On considère une solution $S$ d'acide benzoïque. L'équation de sa réaction avec
l'eau s'écrit~:
    
\hspace{5cm}\ch{C6H5COOH_{(aq)} + H2O_{($\ell$)} <=> C6H5COO^{-}_{(aq)} + H3O^{+}_{(aq)}}

La mesure de sa conductivité a donné la valeur suivante~:
$\sigma=\qty{36.1}{\mS\per\m}$
\begin{enumerate}
    \item Dresser le tableau d'avancement de la réaction.
    \item Donner l'expression de la conductivité $\sigma$ du mélange réactionnel
          en fonction de $\lambda_{\ch{H3O^{+}}}$, $\lambda_{\ch{C6H5COO^{-}}}$,
          du volume $V$ de la solution et l'avancement final $x_{f}$.
    \item Déterminer la valeur de l'avancement final de la dissociation de
          l'acide benzoïque dans l'eau.

          \textbf{On donne~:}
          $\lambda_{\ch{H3O^{+}}}=\qty{35}{\mS\square\m\per\mol}$~;
          $\lambda_{\ch{C6H5COO^{-}}}=\qty{3.23}{\mS\square\m\per\mol}$ et
          $V=\qty{50}{\mL}$.
    \item En déduire les concentrations molaires finales de $\ch{H3O^{+}}$ et
          $\ch{C6H5COO^{-}}$.
    \item Calculer le pH de la solution obtenue.
    \item Déterminer le taux d'avancement final sachant que la concentration de
          la solution est~: $c=\qty{1.18e-2}{\mol\per\L}$
\end{enumerate}
